\documentclass[12pt]{article}
\usepackage[utf8]{inputenc}
\usepackage[T1]{fontenc}
\usepackage[margin=1in]{geometry}
\usepackage{setspace}
\usepackage{xcolor}
\usepackage[colorlinks=true, urlcolor=blue!60!black, linkcolor=blue!60!black]{hyperref}
\setlength{\parindent}{0em}
\setlength{\parskip}{0.45em}
\renewcommand{\baselinestretch}{1.0}
\begin{document}
\begin{center}
{\Large \textit{Yick Wo v. Hopkins}, 118 U.S. 356 (1886)\par}
\vspace{0.3em}
{\normalsize Case Brief --- One-Page Edition\par}
\end{center}

\textbf{Memory Jogger.} The San Francisco laundry case: discriminatory administration of a neutral law violates equal protection (\href{https://supreme.justia.com/cases/federal/us/118/356/}{118 U.S. 356}).

\textbf{Facts.} San Francisco required Board consent to run a laundry in a wooden building, with no standards. Every Chinese applicant, including Yick Wo, whose laundry had passed safety inspections, was refused; nearly every non-Chinese applicant was approved.

\textbf{Procedural History.} Convicted and jailed, Yick Wo was denied habeas by the California Supreme Court.

\textbf{Votes.} Reversed, 9--0. Justice Matthews wrote for the Court.

\textbf{Holding.} The Fourteenth Amendment protects all persons, including noncitizens, and a neutral law administered to discriminate by race or nationality denies equal protection.

\textbf{Reasoning.} The ordinances conferred ``a naked and arbitrary power'' over persons, unlike the uniform curfew in \textit{Barbier}. Holding one's livelihood at another's will is ``the essence of slavery itself,'' and given the admitted disparity, only hostility to the petitioners' race explained it. \textit{Washington v. Davis} later recast the case as proof of purpose.

\textbf{Examples.} \emph{Same-Side:} a standardless board denying nearly all immigrant applicants. \emph{Opposite-Side:} a uniform fire code with disparate impact only. \emph{Fence-Sitter:} a moderate disparity with mixed reasons.

\textbf{Critique.} Its stipulated record makes it an outlier that sets a hard benchmark for proving intent.

\textbf{Key Quotation.} ``Though the law itself be fair on its face and impartial in appearance, yet, if it is applied and administered by public authority with an evil eye and an unequal hand, so as practically to make unjust and illegal discriminations between persons in similar circumstances, material to their rights, the denial of equal justice is still within the prohibition of the Constitution.'' --- Justice Matthews, Opinion of the Court, 118 U.S. at 373--374.

\end{document}
